\section{Desafios práticos}

\desafio{Num projeto Quarkus com um endpoint REST e um bean
\texttt{@ApplicationScoped}, rode o build de produção com
\texttt{-Dquarkus.debug.generated-classes-dir=target/gerado}. Encontre a
classe gerada para o seu bean e leia-a com \texttt{javap -c -p}. Que
instruções aparecem onde você esperaria reflexão?}

\desafio{Ative \texttt{-Dquarkus.debug.print-startup-times=true} e
anote os cinco build steps mais lentos. Remova uma extensão que você não
usa e compare. O que mudou no build e o que mudou no tempo até a
primeira requisição?}

\desafio{Mova uma classe com \texttt{@ApplicationScoped} para um módulo
Maven separado, sem \texttt{beans.xml} e sem índice Jandex. O que
acontece na injeção? Conserte de duas formas diferentes, usando as
fontes do índice listadas na seção~\ref{sec:jandex}, e explique qual
você usaria num time.}

\desafio{Troque, num bean seu, a injeção por campo \texttt{private} por
injeção pelo construtor. Usando a listagem da seção~\ref{sec:gizmo},
explique a um colega por que essa mudança importa mais no executável
nativo do que no modo JVM.}
